% !TEX root = ../main.tex

\chapter{Difficultés rencontrées et solutions apportées}

Au-delà des difficultés techniques déjà détaillées au
Chapitre~\ref{chap:historisation}, plusieurs difficultés d'un autre ordre
ont marqué ce stage.

\section{Montée en compétence sur un existant complexe}

Le projet FPG constituait, au démarrage du stage, une base de code
existante et déjà conséquente. La compréhension de son organisation, de
ses conventions internes et de la logique métier sous-jacente a
nécessité un temps d'adaptation important avant de pouvoir intervenir
efficacement, en particulier sur un sujet aussi transverse que
l'historisation, qui touche de nombreuses entités du système.

\section{Coordination entre deux équipes}

Le travail au sein d'un dispositif impliquant deux équipes distinctes —
le bureau ISI.nc et l'équipe produit FPG — a nécessité une adaptation
constante en termes de communication : reformuler les problématiques
techniques pour un public parfois moins technique côté FPG, et à
l'inverse, faire remonter les contraintes métier auprès de l'équipe
ISI.nc.

\section{Environnement de développement à dépendances externes}

L'environnement de développement reposait sur plusieurs services
externes conteneurisés via Docker (base de données, MinIO) ainsi que sur
un service tiers, JasperReports, pour la génération de rapports PDF. La
gestion de ces dépendances, parfois source de dysfonctionnements
difficiles à diagnostiquer (services non démarrés, versions
incompatibles), a constitué un apprentissage à part entière sur la
robustesse d'un environnement de développement full stack.
